% ===================================================================
%  TABELLA nr. 5 — La chasa e sia construcziun.
%  Traduction 2026 d'apres texto/io, controlee sur texto/fr, texto/ca
%  et texto/oc. Consigne : voir texto/rm/10-tabella-01.tex.
%
%  TROIS COUCHES SUR UN SEUL CHANTIER, ET CHACUNE NOMME UNE ROUTE
%  DIFFERENTE POUR ENTRER DANS LA VALLEE.
%
%  1. LA PIERRE EST LATINE. « murader » le macon, de MURUS ;
%     « cuvrider » le couvreur, de COOPERIRE ; « faver » le forgeron,
%     de FABER ; « vaidrer » le vitrier ; « sarrader » le serrurier ;
%     « scurets » les volets, de OBSCURUS ; « chaltschina » la chaux ;
%     « sablun » le sable ; « crappa » la pierre. Ce sont les metiers
%     qu'on exerce avec ce qu'on a sous la main, et la vallee les
%     nomme depuis Rome.
%
%  2. L'ATELIER EST ALLEMAND. « schrinari » le menuisier vient de
%     Schreiner. Le menuisier, a la difference du macon, ne travaille
%     pas la matiere du lieu : il travaille avec un etabli, des
%     rabots et un tour, c'est-a-dire avec un ATELIER, et l'atelier
%     est venu du nord avec son nom.
%
%  3. LE PAPIER EST ROMAN DU DEHORS. « architect », « impressari »,
%     « plan », « fatschada », « balcun », « terrassa », « veranda »,
%     « mansarda », « salun », « corridor », « vestibul », « alcova ».
%     Tout ce qui se dessine avant d'etre bati, et toutes les pieces
%     que le dessin decoupe, porte un nom venu par le papier.
%
%  LA LIGNE DU TABLEAU 16 LUXEMBOURGEOIS PASSAIT ENTRE CE QUI A UN
%  GUICHET ET CE QUI N'EN A PAS ; ICI ELLE PASSE ENTRE CE QUI SE FAIT
%  A LA MAIN, CE QUI SE FAIT A L'ETABLI ET CE QUI SE FAIT SUR LE
%  PAPIER. Trois etats du meme metier, et chacun a garde le nom qu'il
%  portait quand il est arrive.
%
%  LA NOTE (*) DE L'AUTEUR TOMBE JUSTE AU BON ENDROIT, ET IL FAUT LA
%  LIRE POUR CE QU'ELLE EST. Guignon y manque d'une racine pour la
%  solive et va la chercher en allemand — « Balk-o », de Balken. Il
%  s'en explique : ce n'est ni la poutre ni la petite poutre, et il
%  lui faut un mot technique. Or le romanche, qui a l'allemand
%  derriere le col depuis mille ans, n'a pas pris Balken : il derive
%  de « trav ». CE N'EST DONC PAS LE VOISINAGE QUI FAIT L'EMPRUNT.
%  Le romanche a pris « chegls » et « stelzas » a l'allemand au
%  tableau 2 parce que les jeux lui venaient du nord ; il n'a pas
%  pris Balken parce qu'un charpentier de montagne avait deja son
%  mot. LE CONTACT OFFRE, C'EST LE MANQUE QUI PREND.
%
%  Comme dans texto/io, aucun alinea de ce tableau n'est numerote par
%  l'auteur : la cle est « 01 » pour tout le tableau, et le rang suit
%  les repliques du dialogue.
% ===================================================================

%%K t05-apar-1 apar
\VUsaut{6.0mm}
\VUtitre{15.0pt}{60}{TABELLA nr. 5}
\VUsaut{1.4mm}
\VUtitre{11.4pt}{0}{La chasa e sia construcziun.}
\VUsaut{2.2mm}

%%K t05-tit-1 sub
\VUpk{10.2pt}{40}{Buna scuntrada.}
\VUsaut{3.0mm}

%%K t05-01-1 p
\textsc{Gion}. --- Aug, jau vuless fitg savair co ch'ins constrouescha ina
\VUgras{chasa}.

%%K t05-01-2 p
\textsc{L'aug} (80). --- Quai è fitg simpel, mes ami, vegn cun mai, jau ta mussarai
quella che ser Bernard (79) fa construir e ch'jau vegn ad abitar.

%%K t05-01-3 p
\textsc{Gion}. --- Tge ha pia dissegnà il \VUgras{plan} \textsuperscript{(76)} da questa
chasa?

%%K t05-01-4 p
\textsc{L'aug}. --- L'\VUgras{architect} \textsuperscript{(75)}; el ha preschentà in
dissegn fitg cler ch'è vegnì approvà, e l'\VUgras{impressari}
\textsuperscript{(77)} ha cumenzà las lavurs.

%%K t05-01-5 p
\textsc{Gion}. --- Tge è pia l'impressari?

%%K t05-01-6 p
\textsc{L'aug}. --- Quai è ser Alv, il bab da tes ami Jachen. El dirigia ils
lavurants cun ser Ruf, l'architect.

%%K t05-01-7 p
Ma tge! Qua vegnan gist a temp ser Alv cun sia \VUgras{regla} \textsuperscript{(78)} e ser
Ruf cun ses plan.

%%K t05-01-8 p
Bun di, signurs. Co vai?

%%K t05-01-9 p
\textsc{Ser Ruf}. --- Fitg bain, grazia. Bun di, Gion, tge ch'quest uffant è
creschì!

%%K t05-01-10 p
\textsc{L'aug}. --- Gea, propi, el è in pitschen um. El è fitg serius, ins sto al
infurmar da tut quai ch'el vesa. Oz vuless el savair co ch'ins constrouescha ina
chasa.

%%K t05-tit-2 sub
\VUpk{10.2pt}{40}{Ils lavurants.}
\VUsaut{3.0mm}

%%K t05-01-11 p
\textsc{Ser Alv}. --- Bain, vegn cun mai, amitg, jau ta explitgesch quai
immediatamain, perquai ch'jau hai fitg gugent ils mats che vulan sa instruir.

%%K t05-01-12 p
Ti vesas quel lavurant che tegna ina \VUgras{pala} \textsuperscript{(82)} en maun: quai è in
\VUgras{terraglier} \textsuperscript{(81)}. El chava ina \VUgras{fossa}
\textsuperscript{(46)} per metter il conduct d'aua. El ha era chavà il terren là
nua ch'la chasa vegn a star, per ch'il murader possia auzar ils quatter
\VUgras{mirs}. La part da quels mirs ch'è en terra vegn numnada
\VUgras{fundaments} \textsuperscript{(2)}. Il spazi serrà tranter ils fundaments furma
il \VUgras{chaler} \textsuperscript{(3)}. Quel chaler ha ina pitschna fanestra numnada
\VUgras{fanestrin dal chaler} \textsuperscript{(5)}, ina \VUgras{porta}
\textsuperscript{(4)} ed ina stgala.

%%K t05-01-13 p
Il \VUgras{murader} \textsuperscript{(108)} ha cuntinuà a construir ils mirs, laschond
avertas las giaviadas per las \VUgras{portas} \textsuperscript{(8)} e las
\VUgras{fanestras} \textsuperscript{(17)}.

%%K t05-01-14 p
\textsc{Gion}. --- Ma quel murader na vesai jau betg!

%%K t05-01-15 p
S\textsuperscript{r} \textsc{Alv}. --- Uss ha el finì da lavurar en la chasa. El è
sper la \VUgras{stalla} \textsuperscript{(54)}, sin ses \VUgras{impalcadi}
\textsuperscript{(110)}, el constrouescha il mir da la \VUgras{lavandaria}
\textsuperscript{(56)}.

%%K t05-01-16 p
El ha ina cazzola, in trog ed in \VUgras{plumb} \textsuperscript{(109)}. Ma ti na
vegns betg ad ans algordar da quels nums.

%%K t05-01-17 p
\textsc{Gion}. --- Segir che gea, perquai ch'jau vegn immediatamain a dumandar da
mes aug ina pitschna cazzola ed in trog, e jau sez vegn a far in plumb cun ina
cordetta.

%%K t05-tit-3 sub
\VUsaut{3.0mm}
\VUpk{10.2pt}{40}{Il material.}
\VUsaut{3.0mm}

%%K t05-01-18 p
\textsc{L'aug}. --- Ha! fitg bain. Ma di a mai, Gion, tge ch'ins ha da basegn per
construir ina chasa?

%%K t05-01-19 p
\textsc{Gion}. --- Ins ha da basegn da \VUgras{crappa da taglia} \textsuperscript{(59)} e
da \VUgras{malta} \textsuperscript{(65)}.

%%K t05-01-20 p
\textsc{L'aug}. --- Endretg, guarda quel um cun il grond chapè, sin ses \VUgras{char}
\textsuperscript{(136)}. El è \VUgras{charrader} \textsuperscript{(135)}. Mintga di
maina el qua \VUgras{blocs da crappa} \textsuperscript{(60)}. El descharga ses char en
la curt, ed ils \VUgras{tagliacrappa} \textsuperscript{(83)} taglian ils blocs ed als
resgian cun lur \VUgras{resgia da crappa} \textsuperscript{(85)}. In \VUgras{maunaval}
\textsuperscript{(146)} als porta al murader, che als metta a lieu.

%%K t05-01-21 p
Entant prepara il \VUgras{malteder} \textsuperscript{(87)} la malta. El ha da basegn
da \VUgras{sablun} \textsuperscript{(62)}, da \VUgras{chaltschina}
\textsuperscript{(63)} e d'\VUgras{aua} \textsuperscript{(64)}. La chaltschina è sper el
en in \VUgras{satg} \textsuperscript{(71)}, in \VUgras{servent dal murader}
\textsuperscript{(111)} al maina il sablun sin ina \VUgras{carriola}
\textsuperscript{(112)}, e l'\VUgras{emprendist} \textsuperscript{(113)} è a la pumpa
(58). El emplenescha ina \VUgras{sadella} \textsuperscript{(114)}. La malta è prest
pronta. Il \VUgras{purtader da malta} \textsuperscript{(107)} la prenda e la porta si,
tras la stgala, sin l'impalcadi, nua ch'in murader finescha quella vart da la
chasa.

%%K t05-01-22 p
Ed uss, co vegn numnà il lavurant che lavura il \VUgras{tetg}
\textsuperscript{(31)}?

%%K t05-01-23 p
\textsc{Gion}. --- El è \VUgras{cuvrider} \textsuperscript{(118)}.

%%K t05-tit-4 sub
\VUsaut{3.0mm}
\VUpk{10.2pt}{40}{Il tetg da la chasa.}
\VUsaut{3.0mm}

%%K t05-01-24 p
S\textsuperscript{r} \textsc{Alv}. --- Gea, ma l'emprim vegn il \VUgras{lennari}
\textsuperscript{(115)}.

%%K t05-01-25 p
Il lennari lavura l'\VUgras{ovra da lain} \textsuperscript{(35)}. El assemblescha las
\VUgras{travs} \textsuperscript{(37)} cun \VUgras{chavigls} \textsuperscript{(117)}. Quels
lavurants, là si, cun lur largias chautschas da velur, èn lennaris. In tegna ina
travetta, e l'auter taglia in chavigl cun sia \VUgras{ascha} \textsuperscript{(116)}.
Quel ch'è sper il \VUgras{chamin} \textsuperscript{(41)} è il cuvrider. El inguotta
lattas sin las (*) balkas, ed \VUgras{ardesias} \textsuperscript{(66)} sin las lattas.
Mintgatant metta el tegels.

%%K t05-noto-1 noto
\VUnotes{143.2mm}{%
(*) \VUgras{Balk-o} = D. \textit{Balken}, E. \textit{joist}, F. \textit{solive}, I.
\textit{travicello}, R. \textit{balka}, S. Port. \textit{viga}. Ragisch tecnica
enfin ussa betg adoptada, ma da proponer per translatar il senn dal pled F.
\textit{solive}, ch'è ni ina trav F. \textit{poutre}, ni ina travetta. Quai è ina
staisa da lain quadrada che porta il palantschieu ed il plafun.%
}

%%K t05-01-26 p
Il tetg da la stalla è \VUgras{cuvert cun tegels} \textsuperscript{(55)}, quel da la
chasa è \VUgras{cuvert cun ardesias} \textsuperscript{(32)} e quel da la veranda è
\VUgras{da zinc} \textsuperscript{(24)}.

%%K t05-01-27 p
\textsc{Gion}. --- Lavura il cuvrider era ils tetgs da zinc?

%%K t05-01-28 p
S\textsuperscript{r} \textsc{Alv}. --- Na, ma il \VUgras{lavurant da zinc}
\textsuperscript{(121)} u il \VUgras{lavurant da plumb}, quel che metta la
\VUgras{fanestra da tetg} \textsuperscript{(38)} sin il tetg da la chasa. El metta era
ils \VUgras{tubs da plievgia} \textsuperscript{(43)} ed ils \VUgras{channals dal tetg}
\textsuperscript{(42)}. El suda ensemen il zinc ed il blec. El ha fixà il
\VUgras{parafulmin} \textsuperscript{(40)} e la \VUgras{frizza da vent}
\textsuperscript{(39)} sin il \VUgras{frontispizi} \textsuperscript{(36)}.

%%K t05-01-29 p
\textsc{Gion}. --- Uschia è la chasa finida?

%%K t05-tit-5 sub
\VUsaut{3.0mm}
\VUpk{10.2pt}{40}{Ils utensils.}
\VUsaut{3.0mm}

%%K t05-01-30 p
S\textsuperscript{r} \textsc{Alv}. --- Betg dal tut. Il \VUgras{gipsader}
\textsuperscript{(99)} constrouescha cun \VUgras{matuns} \textsuperscript{(61)} las
paraids che la partan en differentas chombras. El intunga cun \VUgras{gips}
\textsuperscript{(70)} ils \VUgras{plafuns} \textsuperscript{(26)} ed ils mirs. Uss
lavura el il plafun da la \VUgras{veranda} \textsuperscript{(33)}. El ha, sco il
murader, ina \VUgras{cazzola} \textsuperscript{(100)} ed in \VUgras{trog}
\textsuperscript{(101)}.

%%K t05-01-31 p
Lura lavura il schrinari las portas, las fanestras, ils \VUgras{parquets}
\textsuperscript{(16)} ed ils \VUgras{scurets} \textsuperscript{(19)}.

%%K t05-01-32 p
Uss lavura el en la curt sin ses \VUgras{banc da lavur} \textsuperscript{(90)}. El ha
ina \VUgras{resgia} \textsuperscript{(94)} per resgiar il lain (69), ina
\VUgras{plana} \textsuperscript{(91)}, ina \VUgras{morsa} \textsuperscript{(92)}, in
\VUgras{stgalpel} \textsuperscript{(94 bis)}, in \VUgras{cumpass}
\textsuperscript{(95)} ed in \VUgras{martè} da lain \textsuperscript{(95 bis)}.

%%K t05-01-33 p
\textsc{Gion}. --- Ha! ma i è pli divertent da far lain che da murar. Aug, ti vegns
a cumprar per mai in banc da lavur, ina resgia ed ina plana, perquai ch'jau vegn
prest a far ina chamona per Medor.

%%K t05-01-34 p
\textsc{L'aug}. --- Gea, matin.

%%K t05-01-35 p
\textsc{Gion}. --- Tge cuntent ch'jau sun! E tge è quel ch'è en pes sper la porta
d'entrada?

%%K t05-tit-6 sub
\VUsaut{3.0mm}
\VUpk{10.2pt}{40}{Il pictar.}
\VUsaut{3.0mm}

%%K t05-01-36 p
S\textsuperscript{r} \textsc{Alv}. --- El è \VUgras{pictur} \textsuperscript{(102)}. El
stat sin sia stgala cun in vaschè da \VUgras{colur} \textsuperscript{(103)} e ses
\VUgras{penel} \textsuperscript{(104)}. El picta il lain da la porta d'entrada per
ch'el sa mantegnia pli ditg.

%%K t05-01-37 p
El encolla era ils palpiris en las abitaziuns.

%%K t05-01-38 p
Il \VUgras{tapezier} \textsuperscript{(107)}, ch'è sin ina \VUgras{stgala}
\textsuperscript{(106)}, sin il \VUgras{balcun} \textsuperscript{(28)}, metta in
\VUgras{storen} \textsuperscript{(29)}.

%%K t05-01-39 p
Il lavurant ch'è en pes sper la gronda fanestra è il \VUgras{vaidrer}
\textsuperscript{(96)}. El fixescha ils \VUgras{vaiders} \textsuperscript{(18)} cun
\VUgras{mastic} \textsuperscript{(73)}. Quel auter lavurant, mess en schanuglias sper
la porta dal chaler, è \VUgras{sarrader} \textsuperscript{(97)}. El metta ina
\VUgras{sarralia} \textsuperscript{(6)}. Sia \VUgras{lima} \textsuperscript{(98)} è
sper el.

%%K t05-01-40 p
\textsc{Gion}. --- È era sarrader quel lavurant che batta cun ses \VUgras{martè}
\textsuperscript{(84)} sin in toc fier?

%%K t05-01-41 p
S\textsuperscript{r} \textsc{Alv}. --- Pli endretg ditg, el è faver (125). El forgia
\VUgras{ferramenta} \textsuperscript{(67)} per l'\VUgras{electricist}
\textsuperscript{(124)}, ch'è sin il tetg da la \VUgras{remisa}
\textsuperscript{(51)}. El ha ina \VUgras{fusina} \textsuperscript{(126)}, in
\VUgras{ancugn} \textsuperscript{(127)} ed ina \VUgras{tanaglia}
\textsuperscript{(128)}.

%%K t05-01-42 p
L'electricist fixescha il \VUgras{fil} \VUgras{da rom} \textsuperscript{(44)} dal
telefon.

%%K t05-tit-7 sub
\VUpk{10.2pt}{40}{Il culvitar da l'iert.}
\VUsaut{3.0mm}

%%K t05-01-43 p
\textsc{L'aug}. --- En il fund da la \VUgras{curt} \textsuperscript{(45)}, sper il
\VUgras{grigl} \textsuperscript{(47)} ed il \VUgras{portun} \textsuperscript{(48)},
vesas ti il \VUgras{giardinier} \textsuperscript{(129)}. El prepara l'\VUgras{iertin}
\textsuperscript{(49)}. El ha chavà la terra cun sia \VUgras{bade}
\textsuperscript{(131)}, el semna sems e rastella cun il \VUgras{rastè}
\textsuperscript{(130)}. Lura vegn el, cun sia \VUgras{canna d'arusar}
\textsuperscript{(132)}, ad arusar ils \VUgras{quartins} \textsuperscript{(150)}, e las
plantas vegnan a crescher.

%%K t05-01-44 p
Il \VUgras{servent dals chavals} \textsuperscript{(133)}, ch'ha gist tgirà il chaval
ed al ha dà a magliar, nettegia il \VUgras{chalesch} \textsuperscript{(52)} cun ina
spugna, ina \VUgras{pumpa da maun} \textsuperscript{(134)} ed ina sadella d'aua. Lura
al maina el danovamain en la remisa e serra la porta cun il gross \VUgras{sbarr}
\textsuperscript{(53)}.

%%K t05-01-45 p
Uss es ti cuntent, jau crai, ti has gì avunda explicaziuns.

%%K t05-01-46 p
\textsc{Gion}. --- Gea, aug, ma gugent vuless jau vesair il dadens da la chasa.

%%K t05-01-47 p
\textsc{L'aug}. --- Quai capita damaun. Ser Ruf vegn a t'explitgar il plan ed a ta
mussar il num da mintga chombra.

%%K t05-tit-8 sub
\VUsaut{3.0mm}
\VUpk{10.2pt}{40}{L'urden intern da la chasa.}
\VUsaut{2.4mm}

%%K t05-apar-2 apar
{\centering\textit{(Vesair il plan.)}\par}
\VUsaut{2.4mm}

%%K t05-01-48 p
\textsc{L'architect}. --- Vegn, Gion, jau ta fatsch immediatamain explorar las
differentas parts da la chasa.

%%K t05-01-49 p
Entrain nus il \VUgras{plaunterren} \textsuperscript{(7)}. Qua è il buttun da la
\VUgras{campanella} \textsuperscript{(11)}. Nus muntain ils trais \VUgras{stgalims}
\textsuperscript{(14)} dal \VUgras{perrun} \textsuperscript{(9)}, nus passain il
\VUgras{sogl} \textsuperscript{(10)} da la \VUgras{porta d'entrada}
\textsuperscript{(8)} e qua essan nus a pe da la \VUgras{stgala}
\textsuperscript{(13)}.

%%K t05-01-50 p
\textsc{Gion}. --- Quai è il corridor.

%%K t05-01-51 p
\textsc{L'architect}. --- Na, quai è il \VUgras{vestibul} \textsuperscript{(12)}. Il
\VUgras{corridor} \textsuperscript{(a)} è a la fin. A dretga, qua èn il
\VUgras{salun} \textsuperscript{(b)} e la \VUgras{stanza da magliar}
\textsuperscript{(c)}, en fatscha da la stanza da magliar èn la \VUgras{cuschina}
\textsuperscript{(d)} e la \VUgras{stanza dals servents} \textsuperscript{(e)}, lura
avant il \VUgras{biro}\textsuperscript{(f)}. La \VUgras{veranda}
\textsuperscript{(23)} cun sia \VUgras{porta da vaider} \textsuperscript{(25)} è sper
il salun.

%%K t05-01-52 p
Nus muntain immediatamain la stgala. Tegna tai vi dal \VUgras{sustegn da maun}
\textsuperscript{(15)} e dumbra ils stgalims. Els èn quattordesch. Qua è il
\VUgras{sulèr} \textsuperscript{(m)}, nus essan sin l'emprim \VUgras{plaun}
\textsuperscript{(27)}.

%%K t05-tit-9 sub
\VUsaut{3.0mm}
\VUpk{10.2pt}{40}{L'emprim plaun.}
\VUsaut{3.0mm}

%%K t05-01-53 p
En fatscha da la stgala èn il \VUgras{necessari} \textsuperscript{(20)} e la
\VUgras{toaletta} \textsuperscript{(j)}. A dretga ina bella \VUgras{stanza da durmir}
\textsuperscript{(g)} cun duas fanestras vers la \VUgras{fatschada}
\textsuperscript{(1)} da la chasa. A sanestra èn la \VUgras{stanza da bogn}
\textsuperscript{(1)} e la \VUgras{stanza per ils amis} \textsuperscript{(i)} cun ina
gronda \VUgras{alcova} \textsuperscript{(h)}: sper la stanza da durmir è la
\VUgras{stanza dals uffants} \textsuperscript{(k)}. Ella dat sin ina vasta
\VUgras{terrassa} \textsuperscript{(21)}. Sut questa terrassa è in auter necessari
(20) e, vi dal mir, qua è la \VUgras{campana} \textsuperscript{(22)}, che suna per la
tschaina.

%%K t05-01-54 p
\textsc{Gion}. --- Giain nus sin il \VUgras{segund plaun}.

%%K t05-01-55 p
\textsc{L'architect}. --- I n'exista nagin segund plaun. Sut il tetg èn mo
\VUgras{mansardas} \textsuperscript{(33)} ed il granari (34). Nus avain finì
l'exploraziun, nus pudain turnar giu.

%%K t05-01-56 p
\textsc{Gion}. --- Grazia, signur, quai è fitg interessant. Jau vegn immediatamain a
raquintar a mes bab tut quai ch'jau hai vis.

\VUsaut{4.0mm}
\VUfilet{22mm}
